\begin{statement}{E.37}{Proved}{The complete three-adic class-two fibre and the first nonlifting boundary shadow}
Put
\[
 P_{3^n}=B/\bigl(\gamma_3(H)H^{3^n}\bigr),\qquad
 \mathscr J_{3^n}=\operatorname{im}
       \bigl(B\longrightarrow\mathscr D\times P_{3^n}\bigr).
 \tag{E.37.1}
\]
Fix any one of the $432$ multiplier-one source points
of E.24, retaining its literal commutator word.
The complete set of correction images satisfying the
two order relations in $\mathscr J_{3^n}$ has
\[
                         3^{10n+1}
 \tag{E.37.2}
\]
elements.  Its subset lifting to $\mathscr J_{3^{n+1}}$
is exactly its subset lifting through every level,
and has $3^{10n-1}$ elements.
An explicit two-coordinate criterion for this subset
is given in \textup{(E.37.15)} below.

The full inverse-limit fibre of correction images
has the explicit coordinate parametrization
$\mathbb Z_3^{10}$ in \textup{(E.37.13)}.
Every such compatible family is generating in all
the actual common finite sources and is realized
by one profinite commutator word.  It can be prescribed
simultaneously with any of the $2^{10}$ exponent-four
solutions of E.36 and with the whole abelian boundary
tower.

There is an individually specified multiplier-one
commutator shadow in $\mathscr J_3$ with no lift to
$\mathscr J_9$, displayed in
\textup{(E.37.17)--(E.37.18)}.
It still has the abelian-tower compatibility proved
in E.33 and E.36.  The failed projection here is the
additional class-two map $P_9\to P_3$.
\end{statement}
\begin{proof}
We use exact integral doubled-log coordinates
\[
             (v,Z)\in\mathbb Z_3^{10}
                         \oplus\Lambda^2\mathbb Z_3^{10},
       \qquad Z=2z,
 \quad
 (v,Z)(w,T)=(v+w,Z+T+v\wedge w).
 \tag{E.37.3}
\]
Thus the change from E.33 is the actual invertible
map $(v,z)\mapsto(v,2z)$.
Hall collection over $\mathbb Z_3$ gives the
class-two pro-three quotient of $H$, and reduction
modulo $3^n$ gives precisely the geometric kernel
of $P_{3^n}$.  All entries of the based $s,r$
actions are integral in these coordinates.
They are obtained by the same literal substitutions
and signed Schreier rewriting as E.33.
The program verifies their order relations over
$\mathbb Z$ and the exact doubled-log comparison
with the previously saved modulo-three matrices.

We first retain the full source image.
Let $K\subset H^{\rm ab}$ be E.31's actual
eight-column integral kernel lattice.
The forty-five words in
\textup{(E.33.18)--(E.33.19)} have exactly zero
integral $H$-abelianization and form a basis
of the central reduction modulo three.
Their integral class-two central matrix therefore
has determinant a three-adic unit.  Continuous
three-adic powers of these actual words supply
every element of the entire central lattice,
not just its first reduction.
Squaring the words clears the binary source;
two is a unit in $\mathbb Z_3$, so this does not
reduce their three-adic image.
The eight actual $K$-basis words give every
element of $K\otimes\mathbb Z_3$ in abelianization;
the central words then correct their central
values arbitrarily.  Consequently the exact
common source-kernel image is
\[
 \left\{(v,Z):
       v\in K\otimes\mathbb Z_3,\quad
       Z\in\mathbb Z_3^{45}\right\}.
 \tag{E.37.4}
\]
This equality is a statement about the image
with multiplication \textup{(E.37.3)}; no
arithmetic splitting is assumed.
The upper inclusion follows from the integral
abelian image of E.31.  The explicit word
preimages prove the reverse inclusion.
The same argument proves equality after
reduction at every finite level.

The abelian norm equations together with
$Lv=0$ have a primitive integral rank-three
kernel.  A basis parametrization is
\[
 v=B t=
 (-t_1,t_2,-t_1+t_3,t_2,t_3,-t_3,
          t_1-t_3,t_3,-t_1,t_1)^t .
 \tag{E.37.5}
\]
The nonzero Smith factors of these abelian
equations are units at three, so this same
parametrization is exact modulo every $3^n$.
The affine source condition has not been
discarded.  For a vector in $\ker L$, membership
in $K\otimes\mathbb Z_3$ is exactly divisibility
by three of the two functionals
\[
 \begin{split}
 \lambda_1(v)&=v_1-v_4-v_5-11v_9,\\
 \lambda_2(v)&=v_2-11v_4-v_8-v_9.
 \end{split}
 \tag{E.37.6}
\]
They vanish on $K$ and are independent modulo
three on $\ker L$; the index of $K$ there has
three-part $3^2$, proving equality of these
conditions with that lattice.
On $Bt$ the two functionals both reduce to
$t_1-t_2-t_3$.
For each retained source word $f$, the literal
calculation verifies
$\lambda_1([f])=\lambda_2([f])$ modulo three.
Denote this common residue by $c_f\in\{0,1,2\}$.
The complete affine condition is therefore
\[
                     t_1-t_2-t_3\equiv c_f\pmod3 .
 \tag{E.37.7}
\]
The saved rank-three source lattice has determinant
$48$ in the coordinates $t$, so its three-part is
exactly the index three in \textup{(E.37.7)}.
All $432$ residues are saved; their multiplicities
are $90,144,198$ for $c_f=0,1,2$ respectively.

In the doubled-log coordinates the central norm
equations are $MZ+e(t)=0$.
The $90$-by-$45$ integral matrix has a saved
unimodular Smith decomposition
\[
       S_M=U_M M V_M,\qquad
       \operatorname{diag}(S_M)
                     =(1^{35},3,3,0^8).
 \tag{E.37.8}
\]
The certificate checks the full matrices over
$\mathbb Z$, including the unimodular inverses.
Write $Y=V_M^{-1}Z$ and $E(t)=U_Me(t)$.
The nonzero entries before the two factors
three may have the saved unit signs; these
are retained in the formula below.

Define the integral polynomials
\[
 \begin{split}
 Q(t)&=t_1^2+t_2^2+t_3^2+t_1t_2+t_1t_3-t_2t_3,\\
 F(t)&=t_3-t_2+3Q(t),\\
 G(t)&=t_2+t_3-3t_1^2-4t_2^2-2t_3^2
       -3t_1t_2-5t_1t_3+5t_2t_3
       -(t_3-t_2)Q(t).
 \end{split}
 \tag{E.37.9}
\]
Index the Smith coordinates from $0$ to $44$,
as in the certificate.  The exact polynomial
identities are
\[
 \begin{split}
 E_{35}(t)&=(1+t_3-t_2)F(t)+3G(t),\\
 E_{36}(t)&=3(-t_1+t_3),\\
 E_{63}(t)&=2F(t).
 \end{split}
 \tag{E.37.10}
\]
Every other zero-Smith-row polynomial is an
integer multiple of $F(t)$.
These statements are verified coefficient by
coefficient, including every quadratic term.
The row indexed $63$ is among the $90$ equations,
not among the $45$ central unknowns.
Since two is a three-adic unit,
the complete central-solvability obstruction is
\[
                           F(t)=0.
 \tag{E.37.11}
\]
When it holds, the two factors three in
\textup{(E.37.8)} are divisible by construction.
All the central solutions are exactly
\[
 \begin{cases}
 Y_i=-E_i(t)/(S_M)_{ii},&0\le i<35,\\
 Y_{35}=-G(t),\quad Y_{36}=t_1-t_3,\\
 Y_{37},\ldots,Y_{44}\in\mathbb Z_3\ \text{arbitrary},
 \end{cases}
 \qquad Z=V_MY .
 \tag{E.37.12}
\]
Thus no unnamed rational cokernel or divisibility
condition remains in the description.

The equation has an explicit global three-adic
solution graph.  Its derivative with respect
to $t_3$ is $1+3(2t_3+t_1-t_2)$, a unit.
For any $t_1,t_2$, its reduction has the unique
root $t_3=t_2$ modulo three.
Starting with that root, at precision $3^k$
add
$-3^k(F(t)/3^k\bmod3)$ to $t_3$.
This gives the unique compatible roots at every
precision, hence the unique root in $\mathbb Z_3$.
On this graph \textup{(E.37.7)} becomes
$t_1+t_2\equiv c_f$ modulo three.
The promised complete parametrization is
\[
 \begin{gathered}
 u,w,a_{37},\ldots,a_{44}\in\mathbb Z_3,\qquad
 t_2=u,\quad t_1=c_f-u+3w,\\
 t_3=\text{the unique root of }F(t_1,u,t_3)=0,\qquad
 Y_{37},\ldots,Y_{44}=a_{37},\ldots,a_{44},
 \end{gathered}
 \tag{E.37.13}
\]
with the remaining coordinates given by
\textup{(E.37.12)}.
This is a bijective continuous coordinate map
with continuous inverse: recover $u=t_2$,
$w=(t_1+t_2-c_f)/3$, and the eight free Smith
coordinates.  It proves the stated fibre
$\mathbb Z_3^{10}$ with an actual parametrization.

We now compare it exactly with the finite fibres.
Modulo $3^n$, the equation $F(t)=0$ has one
$t_3$ for each $(t_1,t_2)$.
Condition \textup{(E.37.7)} leaves
$3^{2n-1}$ abelian targets.
For each target, \textup{(E.37.8)} gives
an affine central kernel
\[
       (\mathbb Z/3^n)^8\times(\mathbb Z/3)^2.
 \tag{E.37.14}
\]
Its order is $3^{8n+2}$, proving
\textup{(E.37.2)}.
The two extra finite central coordinates are
precisely controlled: a finite solution lifts
to the next level exactly when
\[
          Y_{35}=-G(t),\qquad
          Y_{36}=t_1-t_3\pmod{3^n}.
 \tag{E.37.15}
\]
For necessity, apply
\textup{(E.37.10)} at the next precision and
divide its two Smith equations by three.
For sufficiency, lift $t_1,t_2$ arbitrarily
to $\mathbb Z_3$, take the unique root $t_3$,
and use \textup{(E.37.12)} with arbitrary lifts
of the last eight coordinates.
The first thirty-five coordinates match the
given finite solution by their unit equations,
and the next two match by
\textup{(E.37.15)}.
The affine source residue is unchanged.
This proves lifting through the entire tower,
not just the next level.
It also gives exactly
$3^{2n-1}3^{8n}=3^{10n-1}$ images.

Here is an individual failed lift.  Use E.24's
source point with indexed CM row $0$, ternary
parameters $(0,0)$ and binary parameter $0$;
its residue $c_f$ is zero.
In the boundary quotient $P_3$ prescribe
$v=0$, and in the doubled-log Smith coordinates
prescribe
\[
           Y_{36}=1,\qquad Y_i=0\quad(i\ne36).
 \tag{E.37.16}
\]
Its two torsion coordinates are $(Y_{35},Y_{36})=(0,1)$.
The corresponding BCH central coordinate in the
original E.33 basis is
\[
 \begin{split}
 z={}&e_{14}-e_{15}+e_{16}+e_{19}+e_{1,10}
       -e_{25}-e_{27}+e_{28}-e_{34}\\
    &+e_{36}+e_{37}+e_{38}
       +e_{45}+e_{46}+e_{47}+e_{67}-e_{78},
       \qquad e_{ij}=e_i\wedge e_j,
 \end{split}
 \tag{E.37.17}
\]
with coefficients in $\mathbf F_3$.
It satisfies both order relations: its Smith
error is a multiple of three.
The source fibre is realized by an actual
commutator word.  The certificate expresses
$f^{-1}(0,z)$ in the complete $51$-vector
kernel basis of E.33, retains every basis
relator preimage, and gives the ordered powers
and the collecting commutator corrections.
Replacing those relators by their fourth
powers retains the prescribed modulo-three
values and clears the binary source.
Thus this is a concrete point of the actual
common source, rather than a tuple in an
unrelated coordinate group.

Any commutator-shadow lift to $P_9$ would have
$t_1=t_2=t_3=0$ modulo three, since its abelian
reduction is zero and the basis
\textup{(E.37.5)} is primitive.
Its second divided Smith equation would force
\[
                   Y_{36}=t_1-t_3=0\pmod3,
 \tag{E.37.18}
\]
contrary to \textup{(E.37.16)}.
This excludes the multiplier-one class-two
lift at the specified next quotient.
Its abelian-tower extensions remain exactly
those already constructed in E.33 and E.36.

We finish the exact realizability and generation
comparisons for the surviving full tower.
The affine abelian values in
\textup{(E.37.5)--(E.37.7)} are realized by
the eight source words for $K$, with
three-supported exponents.
The forty-five central preimages proving
\textup{(E.37.4)} then realize the prescribed
central values in \textup{(E.37.12)}.
These are finite products of continuous powers
of actual source words, so give one element
of $\widehat F'$ at every precision.
Use squared words before these powers to
clear the binary source.
The same supported-exponent construction
retains the two-adic and prime-to-six abelian
choices of E.31.
For a prescribed exponent-four solution,
the sixty zero-abelianization words in E.36
give all its remaining central adjustments.
Cube them to kill the ternary source and
use two-supported powers, dividing their
coefficients by the odd unit three.
They act trivially in every three-primary
class-two quotient.  This proves simultaneous
realization with that exponent-four solution
and with every abelian modulus.

Every finite three-primary boundary map just
constructed is an automorphism of $P_{3^n}$.
It preserves $H$ and hence descends to the
characteristic quotient; its action on the
geometric Frattini quotient has rank ten,
as already checked for all the admissible
modulo-three vectors in E.33.
The Frattini criterion gives surjectivity
on the finite geometric three-group and
then on $P_{3^n}$.
For common sources with $\mathscr D$,
the exponent-four quotient and abelian
two-three moduli, use the exact
nilpotent-residual/Frattini argument in
E.36: the boundary factors have genuine
automorphisms, their common $N_B$ image
has a finite lower-central stable term,
and the $\mathscr D$ image kills
$\gamma_5(N_B)$.
Quotienting the resulting pure stable
factor leaves the two rank-$97$
Frattini quotients through $G_6$.
This proves joint generation.
The prime-to-six abelian factors are added
by the same actual CRT and coprime
$H$-image comparison.  No invariance of
the fixed Ferrari source kernel is used.

The exact integral action, Smith matrices,
all polynomial coefficients, every source
affine residue, finite counts and the
individual word certificate are saved in
\texttt{bc\_fermat\_class\_two\_integral\_model.py},
\texttt{bc\_fermat\_class\_two\_3adic\_norms.py},
and \texttt{bc\_fermat\_class\_two\_3adic\_tower.py},
with their executed outputs, JSON and NPZ
certificates.  All end in \texttt{PASS}.
\end{proof}
